% ID: SEL-D0079-OA-E04
% ESTADO: cotejada. Segunda pasada visual y borrador compilado revisados.
% Fuente DOC: Selectividad/sel_2012_matematicas_aplicadas/1 MATEMÁT. APLIC. CC. SS/EXAMEN 1-ANDALUCÍA.doc
% SHA-256 original: a5dfff58dac253bf021b0d2b9409d68987e9ce6e97cfa4abbaf840019a7a18e2
% Página convertida: 1; opción A; ejercicio 4.

La variable ``tiempo de reacción de un conductor ante un obstáculo imprevisto'' sigue una distribución Normal con desviación típica $0.05$ segundos. Al medir dicho tiempo en 50 conductores se ha obtenido un tiempo medio de $0.85$ segundos.
\begin{enumerate}[label=\alph*)]
\item \textbf{(1.25 puntos)} Halle el intervalo de confianza para el tiempo medio de reacción, con un nivel de confianza del $99\%$.
\item \textbf{(1.25 puntos)} \textquestiondown{}De qué tamaño mínimo ha de tomarse una muestra para que el error de estimación no supere $0.01$ segundos, con un nivel de confianza del $95\%$?
\end{enumerate}
